%Poglavlje 1
%
%
\chapter{Teorijska osnova}
\label{chap:theory}
%
U ovom poglavlju izlaže se teorijska osnova neophodna za razumijevanje
implementacije opisane u narednim poglavljima. Prvo se objašnjava format
UART okvira i način na koji se iz njega izvodi brzina prijenosa podataka,
kao i razlog zbog kojeg se prijem vrši na sredini bita. Zatim se opisuju
dijelovi AVR arhitekture koji se koriste u implementaciji: registri opće
namjene, status registar SREG, ulazno-izlazni portovi te trajanje
izvršavanja instrukcija u taktovima.
%
\section{UART protokol}
\label{sec:uart}
%
UART je asinhroni serijski protokol, što znači da predajnik i prijemnik ne
dijele zajedničku liniju takta. Podaci se prenose u okvirima koji se
sastoje od start bita, osam podatkovnih bita i stop bita. Linija je u
mirovanju (\foreignlanguage{english}{idle}) na visokom naponskom nivou, a
start bit se signalizira prelaskom linije na nizak nivo, čime prijemnik
prepoznaje početak novog okvira. Nakon start bita slijedi osam podatkovnih
bita, pri čemu se prvo prenosi bit najmanje težine
(\foreignlanguage{english}{LSB}). Okvir se završava stop bitom, tokom
kojeg se linija vraća na visok nivo, čime se osigurava razmak prije
eventualnog sljedećeg okvira.

Brzina prijenosa (\foreignlanguage{english}{baud rate}) predstavlja broj
bita koji se prenese u jednoj sekundi i definiše se kao recipročna
vrijednost trajanja jednog bita, odnosno $B = 1/T_{bit}$, gdje je
$T_{bit}$ trajanje jednog bita izraženo u sekundama. Pošto predajnik i
prijemnik ne dijele zajednički takt, oba uređaja moraju unaprijed
poznavati istu vrijednost $T_{bit}$ kako bi tajming predaje i prijema bio
usklađen.

Pošto UART protokol ne koristi zajedničku liniju takta, prijemnik mora
sam odrediti trenutak u kojem će očitati vrijednost svakog bita, na
osnovu poznate brzine prijenosa i trenutka detekcije start bita. Ivice
bita, odnosno prelazi između susjednih bita, najosjetljivije su na šum i
na malo odstupanje takta između predajnika i prijemnika, pa bi očitavanje
na samoj ivici moglo dati pogrešnu vrijednost. Očitavanjem na sredini
bita ostavlja se maksimalna vremenska margina sa obje strane, čime se
povećava otpornost prijema na akumulirano odstupanje tajminga tokom
trajanja okvira.
%
\section{AVR arhitektura}
\label{sec:avr}
%
ATmega328P mikrokontroler sadrži 32 8-bitna registra opće namjene, označena od R0 do R31,
organizovana u registar
\foreignlanguage{english}{file} koji je direktno povezan sa
aritmetičko-logičkom jedinicom. Ovakva organizacija omogućava da se
većina aritmetičkih i logičkih operacija izvrši u jednom ciklusu takta,
jer se oba operanda čitaju iz registar \foreignlanguage{english}{file}-a
i rezultat se u istom ciklusu upisuje nazad. Ovi registri su namijenjeni
isključivo za privremeno čuvanje podataka i izvođenje proračuna unutar
CPU jezgra i sami po sebi nemaju nikakav uticaj na fizičke pinove
mikrokontrolera.

Za upravljanje perifernim jedinicama, uključujući i digitalne ulazno-izlazne
pinove, mikrokontroler posjeduje zaseban adresni prostor
ulazno-izlaznih (I/O) registara, fizički odvojen od registar
\foreignlanguage{english}{file}-a registara opće namjene. Dok registri
opće namjene služe isključivo procesoru za privremeno čuvanje podataka,
I/O registri predstavljaju direktno preslikavanje stanja i konfiguracije
hardvera, poput porta, i svaka promjena njihove vrijednosti odmah se
odražava na fizičkom pinu. Ovaj zaseban prostor takođe omogućava da se
pojedinačni bit u I/O registru postavi ili obriše bez potrebe da se
prethodno pročita i ponovo upiše cijeli sadržaj registra, što je posebno
značajno kada je potrebno precizno kontrolisati trajanje operacije u
ciklusima.

Svaki digitalni I/O port, na primjer port D, kontroliše se pomoću tri
odvojena I/O registra. Registar DDRx određuje smjer svakog pojedinačnog
pina tog porta, odnosno da li se pin ponaša kao ulaz ili izlaz. Registar
PORTx ima dvostruku ulogu koja zavisi od stanja odgovarajućeg bita u
DDRx: kada je pin konfigurisan kao izlaz, bit u PORTx direktno određuje
da li je na pinu prisutan nizak ili visok naponski nivo, dok kod
ulaznog pina isti bit uključuje ili isključuje interni pull-up otpornik.
Registar PINx je, za razliku od druga dva, namijenjen isključivo čitanju
i uvijek odražava stvarni trenutni naponski nivo prisutan na fizičkom
pinu, bez obzira na to da li je taj pin konfigurisan kao ulaz ili izlaz.
Na primjer, da bi se neki pin porta D postavio u stanje logičke
jedinice, potrebno je najprije u registru DDRD postaviti odgovarajući
bit na jedan, čime se taj pin konfiguriše kao izlazni, a zatim u
registru PORTD postaviti isti bit na jedan, čime se na tom pinu
uspostavlja visok naponski nivo.

Status registar (SREG) je poseban I/O registar širine 8 bita, u kojem
svaki pojedinačni bit predstavlja zasebnu zastavicu koja nakon izvršenja
većine aritmetičkih i logičkih operacija bilježi neko svojstvo dobijenog
rezultata. Registar sadrži sljedećih osam zastavica:

\begin{itemize}
    \item C (\foreignlanguage{english}{Carry}) - zastavica prijenosa,
        postavlja se kada operacija proizvede prijenos ili posudbu izvan
        granice od 8 bita;
    \item Z (\foreignlanguage{english}{Zero}) - postavlja se kada je
        rezultat posljednje operacije jednak nuli;
    \item N (\foreignlanguage{english}{Negative}) - odražava najznačajniji
        bit rezultata i pokazuje da li je rezultat negativan;
    \item V (\foreignlanguage{english}{Overflow}) - zastavica
        prekoračenja opsega kod aritmetike sa predznakom;
    \item S (\foreignlanguage{english}{Sign}) - stvarni predznak
        rezultata, dobijen kao isključivo ILI zastavica N i V;
    \item H (\foreignlanguage{english}{Half Carry}) - prijenos između
        nižeg i višeg nibla bajta, koristan kod BCD aritmetike;
    \item T (\foreignlanguage{english}{Bit Copy Storage}) - jednobitno
        mjesto za privremeno čuvanje vrijednosti jednog bita prilikom
        prenošenja između registara;
    \item I (\foreignlanguage{english}{Global Interrupt Enable}) -
        globalno omogućava ili onemogućava obradu prekida.
\end{itemize}

Sadržaj ovog registra se može ispitati kako bi program donio odluku o
daljem toku izvršavanja, na primjer da li će se izvršiti uslovni skok,
čime SREG postaje osnovni mehanizam preko kojeg se ostvaruje uslovno
grananje u programu.

Većina AVR instrukcija izvršava se u jednom ciklusu takta, zahvaljujući
jednostepenom pipeline-u u kojem se dohvatanje sljedeće instrukcije
preklapa sa izvršavanjem trenutne. Instrukcije koje mijenjaju tok
programa, poput relativnih skokova i poziva potprograma, instrukcije
koje ispituju bit i uslovno preskaču sljedeću instrukciju, te uslovni
skokovi kada se grananje zaista izvrši, zahtijevaju dodatni ciklus ili
više, jer se u tim slučajevima pipeline mora isprazniti ili se dodatno
pristupa memoriji programa odnosno steku.
%
